\section{A recursion category from a three-point space}\label{sec:cosheaves}

We construct a recursion category $\mathscr R_X$ from cosheaves on a
three-point space. Its restriction splitting contains a consistent
propositional theory with no complete consistent extension represented
by a restriction idempotent.

For a topological space $X$, let $\mathcal O(X)$ be the category
of open sets and inclusions. A \emph{set-valued cosheaf} is a
covariant functor $C:\mathcal O(X)\to\mathbf{Set}$ for which
every open cover $(O_\lambda)_\lambda$ of $O$ gives a coequalizer
\cite[Definition 2.6]{CurryCosheaves}
\[
 \coprod_{\lambda,\mu} C(O_\lambda\cap O_\mu)
 \rightrightarrows \coprod_\lambda C(O_\lambda)
 \longrightarrow C(O).
\]
The two parallel maps are induced by the inclusions into the two
members of the cover; the empty cover requires
$C(\varnothing)=\varnothing$.
The \emph{Sierpi\'nski space} has two points and exactly three open
sets: the empty set, one singleton, and the whole space.

Let $X=\{a,b,c\}$ have opens
$\varnothing,\{c\},\{a,c\},\{b,c\},X$.
The open subspaces $\{a,c\}$ and $\{b,c\}$ are Sierpi\'nski spaces
with common open singleton $\{c\}$. Thus $X$ is obtained by gluing
two copies of the Sierpi\'nski space along that singleton.
Its category of set-valued cosheaves is equivalent to
$\mathbf{Set}^{J}$, where $J$ has objects $0,1,2$ and nonidentity
morphisms $0\to1$ and $0\to2$.
The equivalence sends a cosheaf $C$ to the functor with values
$C(\{c\}),C(\{a,c\}),C(\{b,c\})$ at $0,1,2$, respectively,
and morphisms induced by the inclusions of these open sets.
The cosheaf value $C(X)$ is the pushout of
$C(\{a,c\})\leftarrow C(\{c\})\to C(\{b,c\})$
(Proposition~\ref{prop:fork-cosheaves}).

Let $\mathcal C_{\mathrm{rec}}$ have computably enumerable subsets
of finite powers of $\N$ as objects. A morphism $f:A\to B$ is a
function computed by a partial recursive program defined at every
element of $A$, with values there in $B$. Equality of morphisms is
equality of their functions on $A$.

Put $\mathcal E=\mathcal C_{\mathrm{rec}}^J$.
An object is a covariant functor $Y:J\to\mathcal C_{\mathrm{rec}}$,
given by the objects $Y_j=Y(j)$, for $j=0,1,2$, and the
morphisms $r_i^Y=Y(0\to i):Y_0\to Y_i$, for $i=1,2$.
A morphism $f:Y\to Z$ is a natural transformation: its components
are morphisms $f_j:Y_j\to Z_j$ of $\mathcal C_{\mathrm{rec}}$ satisfying
$r_i^Zf_0=f_i r_i^Y$.
A \emph{c.e.\ subfunctor} $D\subseteq Y$ consists of c.e.\ subsets
$D_j\subseteq Y_j$ such that $r_i^Y(D_0)\subseteq D_i$.
The underlying sets and functions define a functor
$\mathcal E\to\mathbf{Set}^J$. Composing with the quasi-inverse of
the cosheaf equivalence above assigns a set-valued cosheaf to each
object of $\mathcal E$ and a natural transformation to each morphism.

In the notation $\operatorname{Par}(\mathcal C,\operatorname{Mono})$,
$\operatorname{Mono}$ denotes all monomorphisms of the category
$\mathcal C$. Form the partial-map category
\[
 \mathcal A=\operatorname{Par}(\mathcal E,\operatorname{Mono}).
\]
Its morphisms $Y\rightharpoonup Z$ are isomorphism classes of spans
$Y\xleftarrow{m}D\xrightarrow{f}Z$ with $m$ monic; composition
uses pullback. Monomorphisms are componentwise injections.
Thus a partial morphism is equivalently a computably enumerable
subfunctor $D\subseteq Y$ together with a natural transformation
$D\to Z$. Its restriction is the partial identity on $D$, and its range is the
partial identity on the subfunctor of component images.
Proposition~\ref{prop:fork-structure} proves these assertions and
the structural properties used below.

Define an object $V:J\to\mathcal C_{\mathrm{rec}}$ of $\mathcal E$
by the following objects $V_j=V(j)$ and morphisms
$r_i^V=V(0\to i)$:
\begin{equation}\label{eq:universal-fork}
 \begin{gathered}
 V_0=\N^4,\qquad V_1=V_2=\N^2,\\
 r_1^V:V_0\to V_1,\qquad r_1^V(k,z,a,b)=(k,a),\\
 r_2^V:V_0\to V_2,\qquad r_2^V(k,z,a,b)=(k,b).
 \end{gathered}
\end{equation}
We regard $V$ also as an object of $\mathcal A$, whose objects
are those of $\mathcal E$.
For each object $A$ of $\mathcal C_{\mathrm{rec}}$, choose an integer
$m(A)\ge0$ such that $A\subseteq\N^{m(A)}$, with
$m(V_0)=4$ and $m(V_1)=m(V_2)=2$.
Fix recursive numerical encodings of nonempty finite sequences of tuples,
\[
 \operatorname{code}_m:
 \coprod_{n\ge1}(\N^m)^n\xrightarrow{\cong}W_m\subseteq\N
 \qquad(m\ge0),
\]
where each $W_m$ is recursive and the length and entry operations
are recursive. For an object $A$ of $\mathcal C_{\mathrm{rec}}$, put
\[
 A^\#=\{\operatorname{code}_{m(A)}(a_1,\ldots,a_n):
                 n\ge1,\ a_1,\ldots,a_n\in A\}\subseteq W_{m(A)}.
\]
This is a c.e.\ numerical carrier. Let $V^\#$ have components
$(V^\#)_j=(V_j)^\#$.
Its map $(V^\#)(0\to i):(V_0)^\#\to(V_i)^\#$ decodes a sequence,
applies $r_i^V$ to each entry, and encodes the resulting sequence,
for $i=1,2$.
Choose isomorphisms in $\mathcal E$
$V\cong V\times V$, $V\cong V+V$, and $V\cong V^\#$.
Together with their inverses, these form a \emph{frame} at $V$
\cite[Section 3]{Heller}.
Let $\mathscr R_X$ denote the full subcategory of $\mathcal A$
on the objects obtained from $V$ by finitely many applications of
the chosen binary restriction products, binary coproducts, and
nonempty-word formation.
The isomorphisms, constructed in Lemma~\ref{lem:fork-frames}, make
$\mathscr R_X$ isotypical.

Theorem~\ref{thm:fork-turing} constructs a Turing morphism
$\tau:V\times V\rightharpoonup V$ in $\mathcal A$.
For every object $Y$ of $\mathcal A$ and every morphism
$F:Y\times V\rightharpoonup V$ in $\mathcal A$, there is a
total morphism $c:Y\to V$ such that $\tau(c\times1_V)=F$.
By fullness this property holds in $\mathscr R_X$ when
$Y$ belongs to $\mathscr R_X$. Together with the inherited
prodominical structure, it makes $\mathscr R_X$ a recursion category.

Every object $Y$ of $\mathcal E$ admits a total map $m:Y\to V$ in
$\mathcal A$ with a partial inverse $r:V\rightharpoonup Y$
such that $rm=1_Y$ and $mr=\bar r$
(Lemma~\ref{lem:fork-dominance}). In this sense $V$ is a
\emph{universal object} of $\mathcal A$.
These embeddings give the equivalence
$\operatorname{Split}(\mathscr R_X)\simeq\mathcal A$
proved in Proposition~\ref{prop:fork-recursion}.

\subsection{Local connectedness and lifting}

For an object $Y$ of $\mathcal E$ and an element $y\in Y_0$, take
the subfunctor with components $\{y\}$, $\{r_1^Yy\}$, and
$\{r_2^Yy\}$. Its partial identity is a connected restriction
idempotent in $\mathcal A$. The same holds for a subfunctor with
a singleton in component $1$ or $2$ and empty sets in the other
components. These restriction idempotents detect the definedness and value
of every component of a partial morphism, so $\mathcal A$ is locally
connected. Its splitness, weak choice, ranges, coproducts, and
componentwise inverse-image laws are proved in
Proposition~\ref{prop:fork-structure}.
The full inclusion $\mathscr R_X\hookrightarrow\mathcal A$ is a
range functor.
For each complementary pair $c,c'$ on $V$, the section
$s_c:V\to V+V$ of $\nabla_V(c'+c)$ belongs to $\mathscr R_X$
by fullness. Thus this inclusion satisfies the lifting hypothesis of
Theorem~\ref{thm:local-lifting}, which gives an inseparable pair in
$\mathscr R_X$.
The category $\mathscr R_X$ is not restriction-equivalent to the
classical recursion category. Indeed, the terminal object of
$\mathcal E$ has five subobjects, whereas the terminal object of
$\mathcal C_{\mathrm{rec}}$ has two.
A restriction equivalence with the classical recursion category
would extend to the restriction splittings and then restrict to an
equivalence $\mathcal E\simeq\mathcal C_{\mathrm{rec}}$, contradicting
this invariant. Proposition~\ref{prop:fork-distinction} gives the
subobject calculation.

\subsection{Evaluation functors}

Let
$\mathcal K=\operatorname{Par}(\mathcal C_{\mathrm{rec}},\operatorname{Mono})$.
Its objects are c.e.\ numerical carriers, and its morphisms are partial
recursive functions with c.e.\ domains contained in their sources.
The functor
$\mathsf Q_{\mathrm{num}}:\operatorname{Split}(\mathrm{Rec})\to\mathcal K$
constructed in Appendix~\ref{app:fork-structure} is a restriction
equivalence.
For $j\in\{0,1,2\}$, component evaluation defines
\begin{equation}\label{eq:endpoint-functor}
 E_j:\mathcal A\longrightarrow\mathcal K,\qquad
 E_j(Y)=Y_j,\qquad
 E_j(Y\xleftarrow{m}D\xrightarrow{f}Z)
   =(Y_j\xleftarrow{m_j}D_j\xrightarrow{f_j}Z_j).
\end{equation}
For the Turing morphism $\tau:V\times V\rightharpoonup V$ of
Theorem~\ref{thm:fork-turing}, write
$\tau_j=E_j(\tau):V_j\times V_j\rightharpoonup V_j$.

\begin{proposition}\label{prop:endpoint}
The covariant functors $E_j$ preserve restrictions, zeros, ranges,
binary restriction products and coproducts, and finite joins of
restriction idempotents.
For $i=1,2$, the Turing morphism
$\tau:V\times V\rightharpoonup V$ constructed in
Theorem~\ref{thm:fork-turing} is sent
to a Turing morphism on $\N^2$ in $\mathcal K$.
\end{proposition}
\begin{proof}
Component evaluation preserves identities, monomorphisms, and
pullbacks in the total categories. Consequently \eqref{eq:endpoint-functor}
is independent of the representative span and preserves partial-map
composition. The stated operations are represented componentwise by
partial identities, empty sets, images, products, tagged unions, and
finite unions, respectively.

Use the recursive bijection $\operatorname{encode}_2:\N^2\to\N$,
its inverse $\operatorname{decode}_2:\N\to\N^2$, and the acceptable numbering
$(\varphi_p)_{p\in\N}$ chosen in the construction of $\tau$ in
Appendix~\ref{app:fork-structure}. For $i=1,2$, its component has
the formula
\[
 (E_i\tau)((k,p),(\ell,x))
 =\operatorname{decode}_2\varphi_p(\operatorname{encode}_2(\ell,x)).
\]
For a partial recursive map
$F:Z\times\N^2\rightharpoonup\N^2$ with c.e.\ carrier $Z$,
the parameter theorem gives a total recursive map $s:Z\to\N$ such that
\[
 \varphi_{s(z)}
 =\operatorname{encode}_2\circ F(z,-)\circ\operatorname{decode}_2
 :\N\rightharpoonup\N
 \qquad(z\in Z).
\]
Thus $z\mapsto(0,s(z)):Z\to\N^2$ is the required total parameter map.
\end{proof}

For $i\in\{1,2\}$, let $H_i:J\to\mathcal C_{\mathrm{rec}}$
have component $\{*\}$ at $i$ and empty components elsewhere.
For $Y\in\mathcal E$ and $y\in Y_i$, let $p_y:H_i\to Y$ be the total morphism
whose component at $i$ sends $*$ to $y$. Every total morphism
$H_i\to Y$ has this form for a unique $y\in Y_i$. A partial
morphism from $H_i$ is either total or zero, since $H_i$ has just
two restriction idempotents. Hence postcomposition defines a partial
function on these total morphisms, defining
$\Phi_i:\mathcal A\to\Pfn$ by
\[
 \Phi_i(Y)=\operatorname{Tot}(\mathcal A)(H_i,Y),\qquad
 \Phi_i(f)(p)=fp\quad\hbox{when }fp\hbox{ is total}.
\]
If $fp$ is not total, it is zero, and every further composite
remains zero. If it is total, further postcomposition has exactly
the domain specified by the next partial function. This proves
$\Phi_i(gf)=\Phi_i(g)\Phi_i(f)$; identities are immediate.
Let $U:\mathcal K\to\Pfn$ send a c.e.\ numerical carrier to its
underlying set and a partial recursive morphism to the same partial
function. The bijections $p_y\mapsto y$ give a natural isomorphism
$\Phi_i\cong U E_i$.
Indeed, $fp_y=p_{E_i(f)(y)}$ when $E_i(f)(y)$ is defined, and $fp_y=0$
otherwise.
For a restriction idempotent $e$ on $Y$, put
\begin{equation}\label{eq:probe-domain}
 |e|_i:=\{y\in Y_i:ep_y=p_y\}
       =\{y\in Y_i:E_i(e)(y)=y\}.
\end{equation}
This is the c.e.\ subset on which $E_i(e)$ is the partial identity.

The functor $E_0$ does not preserve this Turing morphism's universality.
The total function $F:V_0\to V_0$ given by
$F(k,z,a,b)=(0,0,z,0)$ cannot equal $\tau_0(p,-)$ for any $p\in V_0$.
Let $T\hookrightarrow V\times V$ be the domain subfunctor of
$\tau$ constructed in Appendix~\ref{app:fork-turing}.
If $F=\tau_0(p,-)$, totality of $F$ gives $(p,x)\in T_0$
for every $x\in V_0$. Naturality of $T\to V$ then gives
\[
 r_1^V F(x)=\tau_1(r_1^Vp,r_1^Vx).
\]
The right-hand side is unchanged when only the $z$ coordinate
of $x=(k,z,a,b)$ changes, whereas the left-hand side is $(0,z)$.

\subsection{Propositional theories and component theorem sets}

Let $\zeta:V\to V+V$ be the coproduct isomorphism of
Lemma~\ref{lem:fork-frames}. At each component $j$, the map
$\zeta_j$ selects the summand according to the parity of the
first coordinate in $V_j$. Define complementary restriction idempotents
$q_0,q_1:V\rightharpoonup V$ by
\[
 q_0=\zeta^{-1}(1_V+0_V)\zeta,\qquad
 q_1=\zeta^{-1}(0_V+1_V)\zeta.
\]
Put $d=\tau\Delta_V:V\rightharpoonup V$ and define
$a,b:V\rightharpoonup V$ by
\[
 a=d^*q_0,\qquad b=d^*q_1.
\]
These are the inseparable restriction idempotents supplied by
Theorem~\ref{thm:local-lifting} for the full inclusion
$\mathscr R_X\hookrightarrow\mathcal A$. They remain inseparable
in $\mathcal A$: fullness leaves the restriction idempotents on $V$, their order,
and their complements unchanged.
Let $A,B\subseteq V$ be the c.e.\ subfunctors on which $a,b$
are the partial identities, and write $A_j=A(j)$ and $B_j=B(j)$
for $j=0,1,2$. For $i\in\{1,2\}$, put $d_i=E_i(d)$.
By Proposition~\ref{prop:endpoint},
\begin{align*}
 A_i&=\{x:d_i(x)\text{ is defined with even first coordinate}\},\\
 B_i&=\{x:d_i(x)\text{ is defined with odd first coordinate}\}.
\end{align*}
Let $\mathcal K_2$ be the full subcategory of $\mathcal K$ on the
objects obtained from $\N^2$ by the chosen binary restriction
products and coproducts. Every object of $\mathcal K_2$ is recursively
isomorphic to $\N^2$. The restriction products, coproducts, and
prodominical zero law restrict to this subcategory. By fullness,
$E_i\tau$ and the total parameter maps of
Proposition~\ref{prop:endpoint} belong to $\mathcal K_2$ whenever
their sources do. Thus $\mathcal K_2$ is a recursion category.

The restriction $U|_{\mathcal K_2}:\mathcal K_2\to\Pfn$ preserves
restrictions, zeros, and ranges, since it sends each morphism to its
underlying partial function. It also preserves the chosen products
and coproducts, so it is a range functor.
Complementary c.e.\ domains in an object of $\mathcal K_2$ are recursive,
so their switching sections are total recursive maps and belong to
$\mathcal K_2$ by fullness. The same switching construction was used
in Section~\ref{sec:arithmetic}. With the componentwise coproduct
identification, $E_i\zeta:\N^2\to\N^2+\N^2$ is a coproduct
isomorphism in $\mathcal K_2$.
Theorem~\ref{thm:local-lifting} therefore makes the partial identities
on $A_i,B_i$ inseparable in $\mathcal K_2$. A recursive separator
would define a complemented restriction idempotent on $\N^2$ in
this full subcategory. Hence $A_i,B_i$ are recursively inseparable.

We construct a theory in classical propositional logic from the
disjoint restriction idempotents $a,b$. Define the polynomial functor
$P:\mathcal E\to\mathcal E$ by
$P(Z)=V+Z+(Z\times Z)$ and
$P(f)=1_V+f+(f\times f)$ for morphisms $f$ of $\mathcal E$.
Its initial algebra $\iota:P(S)\to S$ is constructed from finite
formula trees. Initiality means that every morphism
$\alpha:P(Z)\to Z$ of $\mathcal E$ admits a unique morphism $u:S\to Z$
of $\mathcal E$
with $u\iota=\alpha P(u)$.
The three summands give constructors for atoms, negation, and implication.
Write $\theta:V\to S$ for the total atom map
and $\nu:S\to S$ for the total negation constructor.
Let $\mathrm{Prf}$ be the object of finite proofs for the rules below
and $c:\mathrm{Prf}\to S$ its conclusion morphism in $\mathcal E$,
viewed as a total morphism of $\mathcal A$.
Proposition~\ref{prop:fork-proof-initial} defines an endofunctor
$\mathsf H:\mathcal E/S\to\mathcal E/S$ and equips $(\mathrm{Prf},c)$
with an initial algebra structure
$\delta:\mathsf H(\mathrm{Prf},c)\to(\mathrm{Prf},c)$.
The restriction idempotent $t_0=\ran(c)$ represents the theorems.
For $j=0,1,2$, write $S_j=S(j)$, and let $\theta_j$, $\nu_j$,
and $\Rightarrow_j$ be the components of the atom, negation,
and implication morphisms. At component $j$, the logical axiom
schemes, with $\alpha,\beta,\gamma\in S_j$, are
\[
\begin{gathered}
 \alpha\Rightarrow_j(\beta\Rightarrow_j\alpha),\\
 (\alpha\Rightarrow_j(\beta\Rightarrow_j\gamma))
   \Rightarrow_j((\alpha\Rightarrow_j\beta)\Rightarrow_j
                         (\alpha\Rightarrow_j\gamma)),\\
 (\nu_j\beta\Rightarrow_j\nu_j\alpha)
       \Rightarrow_j(\alpha\Rightarrow_j\beta).
\end{gathered}
\]
Modus ponens at component $j$ concludes $\beta$ from $\alpha$
and $\alpha\Rightarrow_j\beta$. The additional axioms are
$\theta_j(x)$ for $x\in A_j$ and $\nu_j\theta_j(y)$ for $y\in B_j$.
For a restriction idempotent $t$ on $S$, denote by $D_t\subseteq S$
the c.e.\ subfunctor on which $t$ is the partial identity.
Theorem~\ref{thm:fork-proof-presentation} constructs this syntax and
proof object, proves that $D_{t_0}$ is the least
subfunctor closed under modus ponens and containing the axioms, and proves
consistency. The one-node proofs give the first two relations below,
and consistency gives the third:
\begin{equation}\label{eq:atomic-inequalities}
 a\le\theta^*t_0,\qquad
 b\le\theta^*\nu^*t_0,\qquad t_0\,\nu^*t_0=0.
\end{equation}
A \emph{theory extending $t_0$} is a restriction idempotent $t\ge t_0$
on $S$ in $\mathcal A$ such that $D_t$ is closed under modus ponens.
It is consistent when $t\,\nu^*t=0$, and complete when
$t\vee\nu^*t=1_S$. The set $(D_{\nu^*t})_j$ consists of the
formulas whose negations belong to $(D_t)_j$.

\begin{theorem}\label{thm:cosheaf-observation}
Fix $i\in\{1,2\}$. For every consistent theory $t$ extending
$t_0=\ran(c)$ on $S$ in $\mathcal A$, define
\[
 P_{i,t}=(E_i\theta)^{-1}(|t|_i),\qquad
 N_{i,t}=(E_i\theta)^{-1}(|\nu^*t|_i).
\]
The restriction idempotents $t,\nu^*t$ are inseparable in $\mathcal A$.
The sets $P_{i,t},N_{i,t}$ are disjoint c.e.\ subsets of $\N^2$ satisfying
$A_i\subseteq P_{i,t}$ and $B_i\subseteq N_{i,t}$.
They are recursively inseparable. The assignment
\[
 t\longmapsto(\omega\times\omega,P_{i,t},N_{i,t})
\]
with the identity of $\omega\times\omega$ for extension inclusions is a
covariant functor from the poset of consistent theories extending
$t_0$ to $\operatorname{Pair}(\mathrm{Rec})$.
Every such theory is incomplete.
\end{theorem}
\begin{proof}
The category $\mathcal A$ has restriction zeros and joins of restriction idempotents
preserved by total inverse image. Its restriction idempotents $a,b$ are inseparable,
and \eqref{eq:atomic-inequalities} gives the hypotheses of
Corollary~\ref{cor:domain-incompleteness}. Hence $t,\nu^*t$ are
inseparable, and $t$ is incomplete.

The atom map $E_i\theta:\N^2\to S_i$ is total recursive.
The component theorem and refutation sets $|t|_i$ and $|\nu^*t|_i$
are c.e.; their inverse images under the total recursive map
$E_i\theta$ are therefore c.e. Consistency and preservation of meets and
zero make them disjoint. Applying $E_i$ to
\eqref{eq:atomic-inequalities} and using $t_0\le t$ gives
the two inclusions. A recursive separator would consequently
separate $A_i,B_i$, which is impossible.

If $t\le t'$, monotonicity gives inclusions of both component
sets, so the identity of $\omega\times\omega$ is the required pair morphism.
The functor laws follow from the identity maps.
\end{proof}

